\chapter{Devices: talking to the bench}\label{ch:devices}

A Device is a connection ATE opens for the schedule: an instrument on the LAN, a serial port,
a datagram peer, a VISA instrument. Four Device Types exist in v1. Each one takes connection
details from the schedule, and the schedule names them --- ATE has no Device discovery.

\begin{center}
\footnotesize
\begin{tabular}{@{}>{\raggedright\arraybackslash}p{3.4cm}>{\raggedright\arraybackslash}p{11.2cm}@{}}
\toprule
Device Type & What it is \\
\midrule
\texttt{DeviceType.TCP} & a raw socket stream: an instrument taking SCPI on port 5025, a UUT
  debug port \\
\texttt{DeviceType.SERIAL} & RS-232 or USB serial, through pyserial \\
\texttt{DeviceType.UDP} & datagrams; one datagram is one message \\
\texttt{DeviceType.VISA} & GPIB, USBTMC and VXI-11 instruments through pyvisa, named by a VISA
  resource string \\
\bottomrule
\end{tabular}
\end{center}

\section{Opening a Device}

\texttt{ATE.open(deviceType, params [, onError])} returns a handle the schedule passes back to
\texttt{read}, \texttt{write}, \texttt{flush} and \texttt{close}. The handle is opaque: treat it
as a value you cannot inspect.

Every field \texttt{open} needs is required --- it has no defaults, and a missing one is an
error. The builders in \texttt{ate.devices} fill in the sensible ones, which is why almost
every schedule uses them:

\begin{center}
\footnotesize
\begin{tabular}{@{}>{\raggedright\arraybackslash}p{1.6cm}>{\raggedright\arraybackslash}p{6.8cm}>{\raggedright\arraybackslash}p{6.2cm}@{}}
\toprule
Device Type & \texttt{open} fields & Builder defaults \\
\midrule
TCP & \texttt{host}, \texttt{port}, \texttt{timeout}, \texttt{readTerm}, \texttt{writeTerm} &
  \texttt{Tcp(host, port)}: timeout 1.0 s, both terminators \texttt{"\textbackslash n"} \\
SERIAL & \texttt{port}, \texttt{baud}, \texttt{dataBits}, \texttt{stopBits}, \texttt{parity},
  \texttt{flowControl}, \texttt{timeout}, \texttt{readTerm}, \texttt{writeTerm} &
  \texttt{Serial(port)}: 9600 baud, 8 data bits, \texttt{Parity.NONE}, 1 stop bit,
  \texttt{FlowControl.NONE}, timeout 1.0 s, both terminators \texttt{"\textbackslash n"} \\
UDP & \texttt{host}, \texttt{port}, \texttt{localPort}, \texttt{timeout} &
  \texttt{Udp(host, port)}: \texttt{localPort} 0 (any), timeout 1.0 s \\
VISA & \texttt{resource}, \texttt{timeout}, \texttt{readTerm}, \texttt{writeTerm} &
  \texttt{Visa(resource)}: timeout 1.0 s, both terminators \texttt{"\textbackslash n"} \\
\bottomrule
\end{tabular}
\end{center}

Times are seconds. \texttt{readTerm} and \texttt{writeTerm} are byte strings; \texttt{""} means
no terminator, which is what a binary protocol needs. UDP has no terminators at all, because a
datagram is already a whole message. \texttt{parity} takes one of \texttt{Parity.NONE},
\texttt{Parity.EVEN}, \texttt{Parity.ODD}, \texttt{Parity.MARK} or \texttt{Parity.SPACE};
\texttt{flowControl} takes \texttt{FlowControl.NONE}, \texttt{FlowControl.RTSCTS},
\texttt{FlowControl.XONXOFF} or \texttt{FlowControl.DSRDTR}.

Each setter sets one field and returns the table, so a params table can be built in one
expression, with no \texttt{:build()} step at the end: \texttt{devices.Tcp(host, port)} then
\texttt{:setTimeout(2.0)} and so on.

\section{Reading, writing, flushing and closing}

For \texttt{read}, \texttt{write} and \texttt{close} every field is optional; a field you leave
out uses what the Device was opened with. \texttt{read} takes \texttt{size}, \texttt{timeout}
and \texttt{readTerm}. \texttt{write} takes \texttt{data} (required), \texttt{writeTerm} and
\texttt{timeout}. \texttt{close} takes nothing. UDP's \texttt{read} and \texttt{write} take
\texttt{size}, \texttt{timeout} and \texttt{data} only.

Data is raw bytes both ways. The terminator is stripped from \texttt{data} and \texttt{count} is
what is left; a \texttt{write}'s \texttt{count} includes the terminator it sent. Convert with
\texttt{string.byte} or \texttt{string.pack} when you need numbers.

From \texttt{open} to \texttt{close} ATE reads every TCP, SERIAL and UDP Device in the
background and buffers what arrives, so a reply that comes before the schedule asks for it is
not lost. \texttt{ATE.read} takes from that buffer first and waits only for what is not there
yet. Each buffer holds \texttt{devices.receiveBufferBytes} (16 MiB by default); on overflow ATE
drops the oldest bytes, logs one WARNING, and the next read result carries
\texttt{overflow = true}. \texttt{ATE.flush(dev)} empties the buffer and returns how many bytes
it discarded --- useful just before sending a command. VISA has no background reader and no
buffer, so \texttt{flush} discards nothing: polling an SCPI instrument with empty reads makes it
log \emph{Query UNTERMINATED} ($-420$) and ties up the bus.

\noindent\textbf{\texttt{tcp.tslua}}\par
\lstinputlisting[language=ATELua]{../examples/test-writer/tcp.tslua}

\section{Every Device Type, one at a time}

UDP is the one with no terminators. Its only partial read is a \texttt{size} that cuts a
datagram short, and the log says so rather than claiming a timeout.

\noindent\textbf{\texttt{udp.tslua}}\par
\lstinputlisting[language=ATELua]{../examples/test-writer/udp.tslua}

A serial port is named by its COM or tty name; pyserial also accepts a URL such as
\texttt{socket://127.0.0.1:5025}, which is how these examples exercise a serial Device without
hardware.

\noindent\textbf{\texttt{serial.tslua}}\par
\lstinputlisting[language=ATELua]{../examples/test-writer/serial.tslua}

VISA addresses an instrument by a resource string. GPIB instruments are
\texttt{GPIB0::5::INSTR}, VXI-11 instruments are
\texttt{TCPIP0::127.0.0.1::5025::SOCKET}. ATE uses whatever pyvisa backend it finds; it does
not ship NI-VISA.

\noindent\textbf{\texttt{visa.tslua}}\par
\lstinputlisting[language=ATELua]{../examples/test-writer/visa.tslua}

\section{Lifetime}

A Device handle lives for the whole Test Run, so a Pre-test Action can open one and every step
can use it through an upvalue. Closing a Device that is already closed does nothing and reports
success. Using a closed Device is a failure, and goes through \texttt{onError}.

ATE closes every Device still open when the run ends, and logs a WARNING naming each one. It is
worth closing what you open, if only to keep the log clean.

\noindent\textbf{\texttt{left\_open.tslua}}\par
\lstinputlisting[language=ATELua]{../examples/test-writer/left_open.tslua}

\cando{open a Device of any of the four Types, read and write it with per-call overrides, and
close it --- and explain what a buffered, overflowing or truncated read means.}
